\begin{proof}
\textbf{(1) Norm expressed by the inner product.}  
By definition,
\[
\|v\|^{2}= \langle v,v\rangle ,\qquad v\in V .
\tag{1}
\]

\textbf{(2) Square the inequality.}  
Since norms are non‑negative, \(\|x+y\|\le\|x\|+\|y\|\) is equivalent to
\[
\|x+y\|^{2}\le(\|x\|+\|y\|)^{2}.
\tag{2}
\]

\textbf{(3) Expand the left‑hand side.}  
Using bilinearity (or sesquilinearity) of the inner product,
\begin{align*}
\|x+y\|^{2}
&=\langle x+y,x+y\rangle \\
&=\langle x,x\rangle+\langle x,y\rangle+\langle y,x\rangle+\langle y,y\rangle \\
&=\|x\|^{2}+2\,\operatorname{Re}\langle x,y\rangle+\|y\|^{2}.
\end{align*}
\tag{3}

\textbf{(4) Expand the right‑hand side and cancel.}  
\[
(\|x\|+\|y\|)^{2}= \|x\|^{2}+2\|x\|\|y\|+\|y\|^{2}.
\]
Substituting this and \eqref{3} into \eqref{2} gives
\[
\|x\|^{2}+2\operatorname{Re}\langle x,y\rangle+\|y\|^{2}
\le
\|x\|^{2}+2\|x\|\|y\|+\|y\|^{2},
\]
hence, after cancelling the equal terms \(\|x\|^{2}\) and \(\|y\|^{2}\),
\[
2\operatorname{Re}\langle x,y\rangle\le 2\|x\|\|y\|
\;\Longleftrightarrow\;
\operatorname{Re}\langle x,y\rangle\le\|x\|\|y\|.
\tag{4}
\]

\textbf{(5) Apply Cauchy–Schwarz.}  
The Cauchy–Schwarz inequality states
\[
|\langle x,y\rangle|\le\|x\|\|y\|.
\tag{CS}
\]
Since \(\operatorname{Re}\langle x,y\rangle\le|\langle x,y\rangle|\), \eqref{CS} implies \eqref{4}.

\textbf{(6) Conclude the squared inequality.}  
Insert the bound \(\operatorname{Re}\langle x,y\rangle\le\|x\|\|y\|\) into \eqref{3}:
\[
\|x+y\|^{2}\le(\|x\|+\|y\|)^{2}.
\]
Both sides are non‑negative; taking the principal square root preserves the inequality:
\[
\|x+y\|\le\|x\|+\|y\|.
\]

\textbf{(7) Degenerate cases.}  
If \(x=0\) or \(y=0\) (or both), the equality \(\|x+y\|=\|x\|+\|y\|\) holds trivially.

\textbf{(8) Equality condition.}  
Equality in the triangle inequality occurs precisely when equality holds in the Cauchy–Schwarz step \eqref{CS}.  Equality in Cauchy–Schwarz is equivalent to linear dependence with a non‑negative proportionality factor, i.e.
\[
\exists\,\lambda\ge0\text{ such that }y=\lambda x\quad(\text{or }x=0\text{ or }y=0).
\]
In this case the vectors point in the same direction and \(\|x+y\|=\|x\|+\|y\|\).

Thus the triangle inequality holds for all vectors in an inner‑product space, with equality exactly in the described linear‑dependence situation. \(\square\)
\end{proof}